\section{Формулы для радиуса сходимости}\label{sec-power-series-02}

\begin{theorem}[Формула Коши--Адамара]\label{thm-cauchy-hadamard}

Для ряда \eqref{eq-power-series}

\begin{equation}
\frac1R=\limsup_{n\to\infty}\sqrt[n]{|c_n|}.
\label{eq-cauchy-hadamard}
\end{equation}

Принимаются соглашения \(1/0=+\infty\) и \(1/(+\infty)=0\).

\end{theorem}

\begin{proof}

Сдвигом полагаем \(z_0=0\) и обозначаем \(L=\limsup_n\sqrt[n]{|c_n|}\).

\textbf{Случай \(L=0\).} Для каждого фиксированного \(z\) имеем \[
\sqrt[n]{|c_nz^n|}=|z|\sqrt[n]{|c_n|}\longrightarrow0.
\] По коренному признаку Коши ряд сходится абсолютно при всех \(z\), то есть \(R=+\infty\).

\textbf{Случай \(L=+\infty\).} Для каждого \(z\ne0\) существует подпоследовательность \(n_k\), для которой \(|z|\sqrt[n_k]{|c_{n_k}|}\to+\infty\). Тогда члены \(c_{n_k}z^{n_k}\) не стремятся к нулю. Ряд расходится при \(z\ne0\), а при \(z=0\) равен \(c_0\). Значит, \(R=0\).

\textbf{Случай \(0<L<+\infty\).} Если \(|z|<1/L\), выбираем \(\eta>0\) так, чтобы \(|z|(L+\eta)<1\). По определению верхнего предела для всех достаточно больших \(n\), \[
\sqrt[n]{|c_n|}<L+\eta,
\qquad\sqrt[n]{|c_nz^n|}<|z|(L+\eta)<1.
\] По признаку Коши ряд сходится абсолютно.

Если \(|z|>1/L\), выбираем \(0<\eta<L\) с \(|z|(L-\eta)>1\). По определению верхнего предела существует бесконечная подпоследовательность \(n_k\) с \(\sqrt[n_k]{|c_{n_k}|}>L-\eta\). Вдоль неё \[
\sqrt[n_k]{|c_{n_k}z^{n_k}|}>|z|(L-\eta)>1,
\] так что необходимое условие сходимости нарушено. Следовательно, \(R=1/L\).

\end{proof}

\begin{cor}[Частные формулы]\label{cor-radius-limits}

Если существует предел \(\lim_n\sqrt[n]{|c_n|}\), то его можно использовать вместо верхнего предела в формуле \eqref{eq-cauchy-hadamard}.

Если коэффициенты \(c_n\) ненулевые начиная с некоторого номера и существует (в том числе бесконечный) предел отношений, то

\begin{equation}
\frac1R=\lim_{n\to\infty}\frac{|c_{n+1}|}{|c_n|}.
\label{eq-radius-ratio}
\end{equation}

\end{cor}

Условие о ненулевых коэффициентах нужно, чтобы отношения были определены. Ряды с пропущенными степенями следует проверять непосредственно или по формуле \eqref{eq-cauchy-hadamard}.
